\documentclass{article}
\usepackage{graphicx} % Required for inserting images
\usepackage{amsmath}
  \RequirePackage[french]{babel}
  \RequirePackage[T1]{fontenc}
  \RequirePackage[utf8]{inputenc}

\title{STT-2902 - Modélisation statistique}
\author{Julien Miron}


\begin{document}

\section*{Équipe 4 – Tests d'hypothèses}

\textbf{Session :} Automne 2026

\textbf{Style imposé : procès simulé / débat.} Votre présentation doit prendre la forme d'un procès où l'hypothèse nulle est l'accusée : une partie de l'équipe défend $H_0$ (présomption d'innocence), l'autre présente la preuve (les données) contre elle, et l'auditoire joue le rôle du jury. Ce style est obligatoire.

\vspace{0.5cm}
\textbf{Objectif :} Expliquer la logique des tests d'hypothèses (hypothèses nulle et alternative, seuil, valeur-p, erreurs de type I et II) à travers l'analogie du procès.

\vspace{0.5cm}
\textbf{Déroulement imposé} :
\begin{itemize}
    \item Une mise en contexte présentant l'affaire (le scénario statistique à tester) et les rôles de chacun;
    \item Le procès : présentation de la preuve (données et statistique de test), plaidoiries des deux parties, délibération du jury (l'auditoire);
    \item Un verdict rendu par le jury, suivi de la comparaison avec la conclusion formelle du test statistique;
    \item Un retour pédagogique reliant chaque élément du procès au vocabulaire statistique.
\end{itemize}

\vspace{0.5cm}
\textbf{Requis} :
\begin{enumerate}
    \item Utiliser un scénario concret et réaliste (ex. : la moyenne d'un groupe est-elle différente d'une valeur cible?);
    \item Clairement expliquer la valeur-p et son rôle dans la décision;
    \item Illustrer graphiquement les erreurs de type I et II pour votre test, et les relier aux verdicts possibles du procès (condamner un innocent, acquitter un coupable);
    \item Faire voter le jury (l'auditoire) et comparer sa décision avec la conclusion du test au seuil $\alpha = 0{,}05$.
\end{enumerate}

\end{document}
